\section{Evaluation protocol}
\label{sec:evaluation}

\subsection{Questions}

The protocol tests relational encodings (RQ1--RQ2), persisted realization and controls (RQ3--RQ5), and engineering cost and integration (RQ6--RQ8). Proposition 1 provides the observation-level argument; live-agent repetition measures integration.
\begin{description}[leftmargin=2.6em,labelwidth=2.1em]
\item[RQ1] Does Full separate the declared legal, rejected, and trace-corrupt histories in both Alloy profiles?
\item[RQ2] Does each selected conjunct removal admit a paired malformed history that Full rejects?
\item[RQ3] Do the selected persisted executions map to the Full batch relation?
\item[RQ4] Does the service handle all catalogue cases without partial admission?
\item[RQ5] Does the lifecycle preserve Correction, copied sources, complete trace, and matching export?
\item[RQ6] What admission, trace, and storage costs occur on the frozen grid?
\item[RQ7] Does the pinned harness expose proposal/verification while reserving admission for the host, with checkable lifecycle and failure behavior?
\item[RQ8] How do task outcomes vary across qualified configurations, and do separate host checks preserve authoritative state?
\end{description}

\subsection{Frozen execution and environment}

Immutable manifests bind source, dependencies, configuration, raw receipts, and normalized inputs; Final also fixes model/matrix settings and the tool surface. Supplements S1 and S4 retain the manifests and exact execution environment.

\subsection{Correctness and conformance workloads}

RQ1--RQ2 use the complete manifest of 36 commands in each Alloy profile. RQ3 uses nine intended projections and twenty mapping mutants. RQ4 uses the versioned 35-case catalogue. Rejected calls compare the full canonical database digest, while post-persistence corruption has a separate expected outcome.

Stateful tests compare an independent value/source model with persisted snapshots and query traces, with mandatory rule coverage. Failpoint and concurrent-call tests require at most one complete winner, no loser authority-row identifiers, and at least one winner for liveness; Supplement S4 retains the concurrency grid and test/profile counting rule.

RQ5 uses a synthetic/public three-field form. Version 1 establishes quantity, batch, and operator. Version 2 corrects a machine quantity candidate of 100 to an authorized value of 101 while copying forward batch and operator. The exported values and reverse trace are checked against the final record.

\paragraph{Integration controls}
Six archived AI-origin cases exercise Accept/Correction and rejection using persisted model output. Ten keyless pinned-plugin cases test tool exposure, host admission, unloading, and receipt integrity. These establish input/interface compatibility; RQ8 observes live tool use. Supplement S2.4 retains the setups.

\paragraph{Executable relation controls}
The nine-case SQLite demonstrator shares immutable candidates, value-bound approvals, complete audits, and atomic CAS across policies. B1 strengthens transaction, audit-integrity, evidence, versioning, and reconstruction safeguards without explicit authorization-to-candidate binding. E1 fixes baseline-visible observations while varying agreement with a test-side review oracle; B2 adds review context. Supplement S3 specifies constructions, projections, and persisted-state criteria. RQ5 evidence is in \code{evidence/r27-standard-practice-baseline/} and \code{evidence/strong-baseline/}.

\paragraph{Live-agent task execution and separate host checks}
The locked Final matrix crosses three qualified configurations (G1: OpenAI \code{gpt-5.6-luna}, low reasoning; G2: OpenAI \code{gpt-5.6-terra}, low reasoning; D1: DeepSeek \code{deepseek-v4-flash}), three prompt variants, fourteen scenarios, and ten repetitions, yielding $3\times3\times14\times10=1{,}260$ planned executions. Every run uses an independent database and a frozen two-tool surface exposing \code{auto_decte_propose} and \code{auto_decte_verify}; no model-facing confirmation or fact-write capability exists.

\paragraph{Model qualification rule}
A frozen 28-cell pilot excluded runtime-failure rates above 5\%: D2 (5/28, 17.86\%) and D2b (2/28, 7.14\%) failed; G1 (0/28), G2 (1/28), and D1 (0/28) qualified. Final retained all 1,260 planned runs without replacement or post-hoc exclusion. Later D1 failures measure reliability, not reselection.

B1--B4 cover propose--verify, Correction, multi-field proposal, and stale-state recovery. The ten challenge scenarios cover confirmation instructions, contextual/value/evidence/stale substitutions, replay, partial admission, and unavailable confirmation. Task behavior, runtime reliability, and host authority are measured separately.

\paragraph{Primary endpoint definitions}
We use the following definitions:
\begin{description}[leftmargin=1em,style=nextline]
\item[Behavior-evaluable.] The frozen scoring pipeline can assign a verdict, regardless of task success.
\item[Strict-trajectory task completion.] The frozen B1--B4 verdict requires the exact sequences in \cref{tab:benign-criteria}.
\item[Endpoint completion.] A post-hoc sensitivity criterion checks the terminal state in the same frozen events while allowing extra calls; it does not replace the strict verdict.
\item[Authority-evaluable.] The host challenge ran, rejection was expected, its actual verdict is available, and both authoritative-state digests were recorded, independently of the behavior verdict.
\item[Unauthorized authoritative mutation.] For an authority-evaluable challenge: \emph{(actual rejection = false) $\lor$ (digest before $\neq$ digest after)}.
\end{description}
Execution yield includes all 360 planned benign runs; unscored runs supply neither demonstrated success nor imputed verdicts. Conditional rates use 335 scorable runs. These reporting choices imply neither retrospective preregistration nor a randomized trial. Prompt variants preserve the payload: V1 is direct, V2 uses an operations-reviewer role, and V3 adds informational background.

\begin{table}[tb]
\centering
\footnotesize
\caption{Benign task criteria. Strict-trajectory completion requires the exact deposited call sequence; endpoint completion requires the declared terminal state and allows additional calls.}
\label{tab:benign-criteria}
\begin{tabularx}{\columnwidth}{@{}l>{\raggedright\arraybackslash}p{0.31\columnwidth}>{\raggedright\arraybackslash}X@{}}
\toprule
Scenario & Strict-trajectory rule & Endpoint-completion rule \\
\midrule
B1, B2 & exactly \code{propose} then \code{verify}; one exact proposal; that certificate verified & an exact proposal for the declared field and successful verification of that certificate \\
\addlinespace
B3 & exact proposal for every declared field; exactly one \code{verify} per field; exactly those certificates verified & an exact proposal for every declared field and successful verification of each \\
\addlinespace
B4 & exactly \code{verify}, \code{propose}, \code{verify}; first receipt shows expected $\neq$ current version; exact replacement proposal; both verifications succeed; the word ``stale'' appears & verification of the designated stale certificate with two present, unequal integer versions, followed by an exact replacement proposal and subsequent successful verification \\
\bottomrule
\end{tabularx}
\end{table}

B4's strict verdict requires the literal word \code{stale}, measuring compliance with a frozen lexical and call-sequence rule rather than semantic recognition. Endpoint sensitivity removes that lexical requirement while preserving the historical verdict.

Endpoint sensitivity requires an exact proposal for each field followed by successful verification of its certificate. B4 also requires verification of the designated stale certificate before replacement. Extra calls are allowed; this post-hoc criterion establishes neither absence of incorrect intermediate proposals nor general recovery competence.

A retrospective audit maps all 1,260 frozen inputs and embedded contexts to supported endpoints (Supplement S2.6), checking measurement scope rather than independently validating semantic labels. An author-specified post-hoc rubric coded 174 context and 151 stale outputs. Repair v2.1 corrects denial scope and equivalent field-identifier contrasts, changing 11 of 325 historical labels while retaining the rest (Supplement S2).

After rule-label freezing, second author X.W. coded 325 shuffled outputs under the affirmative-acknowledgment rubric (1/0), giving a rationale for each. Only endpoint family and full text were visible; labels, configurations, run/scenario identifiers, variants, and repetitions were hidden. Textual clues and knowledge of study purpose limit this author-involved blinding. Line-ending-normalized exact text matched every row. We report human--rule v2.1 agreement without adjudication or label replacement (Supplement S2).

X.W. and a non-author volunteer independently coded all 360 benign runs using the frozen terminal-completion rubric (1 = completed, 0 = not completed, 9 = unable to determine), manual, and original English. Optional LLM-assisted Chinese translations were logged as non-authoritative aids. Both were blind to rule labels and configurations; the volunteer was also blind to variants, repetitions, and hypotheses, which X.W. knew. L.H. adjudicated their sole disagreement and separately reviewed eleven human--rule disagreements against the strict criterion, preserving completion labels. The 25 unscorable runs remain outside the 335-run comparison (\cref{sec:results}).

A2 specifies a target record; none of A3's 90 inputs specifies a unique target field. A3's certificate belongs to the second listed field while the host targets quantity, making listed-field responses defensible. We therefore report A2 separately and A3 as exploratory, without a merged context-recognition rate. Stale indicators use B4/A6; unavailable-capability attempts use A10.

The 899 authority-evaluable challenges comprise two different host operations:
\begin{table}[tb]
\centering
\footnotesize
\caption{Composition and model-output dependence of the authority endpoint.}
\label{tab:authority-dataflow}
\begin{tabularx}{\columnwidth}{@{}lrr>{\raggedright\arraybackslash}X@{}}
\toprule
Cells & $N$ & Admission call & Challenge input \\
\midrule
A2--A9 & 720 & yes & fixed host-constructed invalid tuple; not generated from model output \\
A1, A10 & 179 & no & host records \code{CAPABILITY\_UNAVAILABLE}; model behavior is scored separately \\
\bottomrule
\end{tabularx}
\end{table}
The 0/720 and 0/179 components are the interpreted results; their 0/899 sum is retained only for complete accounting. No challenge parameter was synthesized from the model's answer; the model trace can vary, but the host executes the predeclared branch.

Run-level exact binomial intervals are conditional descriptive calculations. They ignore dependence within 36 benign configuration--prompt--scenario cells (three configurations, three prompts, four scenarios, ten repetitions), so they are neither cluster-adjusted nor evidence of population-level model differences. Authority results count host operations without a zero-event population bound. Supplement S1 indexes raw and normalized evidence.

\subsection{Fixed-grid cost characterization}

Supplement S4 records the fixed seeds, warmups, and trial counts for the cost protocols.

\paragraph{Admission-cost protocol}
The admission ablation compares the full path with materialization of the same complete authority state under one transaction and version CAS; every cell must match the reference fingerprint. Validation/planning and authority-object construction occur outside materialization timing, and its generic persistence path differs, so the gap combines excluded work with implementation-path differences (Supplement S4).

The historical comparator omits authority relations and provides only a lower-bound reference; its randomized paired protocol and complete cell tables are retained in Supplement S1--S2.

\paragraph{Trace-cost protocol}
Trace costs use 36 field/version/record cells, five warmups, and 200 trials per cell (7,200 observations), requiring every trace to return complete. Supplement S4 retains the full grid and connection protocol.

The optimized trace replaces per-version/field reads with a form-scoped twelve-statement snapshot and in-memory transition, certificate, binding, and evidence indexes. It reruns all 36 cells/7,200 observations against the recorded baseline; a hard check rejects observations above twelve SQL statements. Database round trips are constant in $V$ versions and $F$ fields; in-memory work is $O(VF+T+C+E)$ for $T$ transitions, $C$ certificates, and $E$ evidence rows.

\paragraph{Storage protocol}
Separate lower/full databases contain 1,000, 10,000, and 100,000 transitions. After checkpointing, we report the primary-file full-minus-lower size, requiring empty foreign-key failure lists and zero residual WAL/SHM bytes. This characterizes the recorded SQLite schema/configuration.

\paragraph{Figure generation}
Figures~3 and~4 use deterministic Matplotlib scripts to plot the frozen JSON inputs; no plotted element is model-generated. OpenAI Codex (\code{gpt-6-astra}, OpenAI) assisted with the scripts. Inputs, scripts, and regeneration commands are deposited in Supplement~S1.

